\documentstyle[%


righttag%


,11pt%


,amssymb%


,russian%               remove in Eng.


,izvru%                 Set izven% in Eng.


]{amsart}





%





\newcommand{\gea}{\geqslant}


\newcommand{\lea}{\leqslant}


\newcommand{\vp}{\varphi}


\newcommand{\vk}{\varkappa}


\newcommand{\ve}{\varepsilon}


\newcommand{\N}{\Bbb{N}}


\newcommand{\R}{\Bbb{R}}


\newcommand{\bH}{\Bbb{H}}


\newcommand{\rL}{\cal{L}}


\newcommand{\intl}{\int\limits}


\newcommand{\suml}{\sum\limits}


\newcommand{\sgn}{\operatorname{sgn}}


\newcommand{\ol}{\overline}


\newcommand{\la}{\langle}


\newcommand{\ra}{\rangle}


\numberwithin{equation}{section}


\newcommand{\tts}[1]{}





\theoremstyle{definition}


\theorembodyfont{\rm}


\newtheorem{cornonum}{\hskip\parindent Следствие}


\renewcommand{\thecornonum}{}


\renewcommand{\baselinestretch}{0.98}





%\hoffset=-15mm%                  For Eng.


\setcounter{page}{38}


\god{2005}


\nomer{4~(515)} %                        Remove in Eng.


%\nomer{Vol.\,49, No.\,4}%               Set in Eng.


%\str{\pageref{begin}--\pageref{end}}%  Set in Eng.


\udk{517.968:519.6}





\author{\it Л.Б.\,Ермолаева}


% NB:   ^^ remove in Eng.


\title{Решение обыкновенных интегродифференциальных


уравнений методом осциллирующих функций}





\date{}


%\pagestyle{myheadings}%         Set in Eng.


\pagestyle{plain} %              Remove in Eng.


\begin{document}


%\thispagestyle{plain}%          Set in Eng.


\maketitle


\label{begin}








\section*{Введение}





Данная статья, являющаяся естественным продолжением


работ [1]--[3], посвящена теоретическому обоснованию в смысле


([4], гл.\,14; [5], гл.\,1) метода осциллирующих функций (подобластей) 


для обыкновенных интегродифференциальных уравнений и некоторых их 


обобщений. В частности, доказывается сходимость метода при минимальных


(к настоящему времени) предположениях относительно исходных данных и 


устанавливаются эффективные (в том числе неулучшаемые по порядку) оценки 


погрешности в зависимости от структурных свойств исходных данных. 


При этом существенным образом используются как результаты, так и 


обозначения из [2], [3].








\section{Постановка задачи}





Рассматривается линейное операторное уравнение


\begin{equation}


Ax\equiv x^{(m)}(t)+B(x;t) = y(t),\quad -1\lea t\lea 1,


\end{equation}


при краевых условиях


\begin{equation}


R_l(x)=0,\quad l=\ol{0,m-1};


\end{equation}


здесь $B$ --- линейный (в том числе интегродифференциальный) оператор с 


областью значений в пространстве $L_1(-1,1)$, $y\in L_1(-1,1)$,


а $R_l$ --- линейно независимые функционалы в пространстве


$C^{m-1}[-1,1]$, $m$ --- целое неотрицательное число, причем условия (1.2) 


при $m=0$ отсутствуют.





Приближенное решение задачи (1.1)--(1.2) ищется в виде многочлена


\begin{equation}


x_n(t)=\suml_{k=1}^{n+m}\alpha_k t^{k-1},\quad t\in [-1,1],\quad n\in \N,


\end{equation}


коэффициенты $\alpha_k=\alpha_{k,n}$ которого определяются по методу 


подобластей из системы линейных алгебраических уравнений (СЛАУ)


\begin{gather}


\suml_{k=1}^{n+m}\alpha_k \int_{t_{i-1}}^{t_i} A(t^{k-1};t)dt =


\int_{t_{i-1}}^{t_i}y(t)dt,\quad i=\ol{1,n},\\


%


\suml_{k=1}^n\alpha_k R_l(t^{k-1})=0,\quad l=\ol{0,m-1};


\end{gather}


здесь $\{t_k=t_{k,n}\}_0^n$ --- некоторая система узлов из $[-1,1]$, 


выбор которых имеет, как будет показано ниже, существенное значение


для обоснования сходимости и оценки погрешности метода.





Исследование метода (1.1)--(1.5) будем проводить в парах функциональных 


пространств $(W_2^m(\rho); L_2(\rho))$ и $(C^m,C)$, где


$C=C[-1,1]$, $L_2(\rho)=L_2(\rho;[-1,1])$ --- хорошо известные пространства 


с обычными нормами, а $C^m = C^{(m)}[-1,1]$ --- пространство $m$~раз


непрерывно дифференцируемых на $[-1,1]$ функций с нормой


$$


\|x\|_{C^m}=\suml_{i=0}^m\|x^{(i)}(t)\|_C,\quad x\in C^m;


$$


$W_2^m(\rho) = W_2^{(m)}(\rho;[-1,1])=\{x(t)\in C^{m-1}:


\exists\, x^{(m)}(t)\in L_2(\rho)\}$


--- весовое пространство Соболева с нормой


$$


\|x\|_{W_2^m(\rho)} =


\suml_{i=0}^m \|x^{(i)}(t)\|_{L_2(\rho)},\quad x\in W_2^m(\rho),


$$


причем выбор весовой функции $\rho=\rho(t)=(1-t^2)^{1/2}$ продиктован, 


как и в [1], [2], свойствами используемых ниже многочленов


Чебышева $T_n(t) = \cos n \arccos t$ ($n+1\in \N$, $t\in [-1,1]$)


и их производных.








\section{Основные результаты}





Ниже за узлы $t_k = t_{k,n}$ ($k=\ol{0,n}$) возьмем корни многочлена 


Чебышева $T_{n+1}(t)$ и соответственно экстремальные точки многочлена


Чебышева $T_n(t)$:


\begin{alignat}{2}


t_k &= \cos \frac{2k+1}{2n+2}\pi,&\quad k&=\ol{0,n},\quad n\in \N;\\


%


t_k &= \cos\frac{k\pi}{n},&\quad k&=\ol{0,n},\quad n\in\N.


\end{alignat}





Для вычислительной схемы (1.1)--(1.5), (2.1)--(2.2) справедливы 


следующие теоремы.





\begin{thm}


Пусть выполнены условия


\begin{itemize}


\item[а)] оператор $B:W_2^m(\rho)\to L_2(\rho)$ вполне непрерывен$;$


\item[б)] уравнение $x^{(m)}(t)=0$ при условиях $(1.2)$ имеет лишь 


тривиальное решение$;$


\item[в)] задача $(1.1)$--$(1.2)$ имеет единственное решение


$x^*(t)\in W_2^m(\rho)$ при любой правой части $y(t)\in L_2(\rho)$.


\end{itemize}





Тогда при всех $n\gea n_0$ $(n_0\in \N$ определяется свойствами


оператора $B)$ СЛАУ $(1.4)$--$(1.5)$, $(2.1)$--$(2.2)$ имеет единственное 


решение $\alpha_1^*,\alpha_2^*,\dotsc, \alpha_{n+m}^*$. Приближенные решения


\begin{equation}


x_n^*(t)=\suml_{k=1}^{n+m}\alpha_k^*t^{k-1},\quad -1\lea t\lea 1,


\tag{1.3${}^*$}


\end{equation}


сходятся при $n\to \infty$ к точному решению $x^*(t)$ со скоростью, 


определяемой соотношениями


\begin{equation}


E_{n-1}({x^*}^{(m)})_{L_2(\rho)}\asymp \|x^*-x^*_n\|_{W_2^m(\rho)} \asymp


E_{n+m-1}(x^*)_{W_2^m(\rho)},


\end{equation}


где символ $\asymp$ означает знак слабой эквивалентности.


\end{thm}





Здесь и далее $E_{n+m-1}(\vp)_{W_2^m(\rho)}$ --- наилучшее приближение 


функции $\vp\in W_2^m(\rho)$ всевозможными алгебраическими многочленами


вида (1.3) в метрике пространства $W_2^m(\rho)$,


а $E_{n-1}(f)_{L_2(\rho)}$ --- наилучшее приближение в пространстве


$L_2(\rho)$ функции $f\in L_2(\rho)$ алгебраическими многочленами степени 


не выше $n-1$.


\medskip








{\bf Доказательство.}


Положим $X=\{x\in W_2^m(\rho):R_l(x)=0$, $l=\ol{0,m-1}\}$


и $Y=L_2(\rho)$ с нормами соответственно


$$


\|x\|_X = \suml_{i=0}^m\|x^{(i)}(t)\|_Y,\quad x\in X;\quad


\|y\|_Y = \bigg\{\int_{-1}^{+1} \rho(t)|y(t)|^2 dt\bigg\}^{1/2},\quad y\in Y.


$$


Тогда задача (1.1)--(1.2) может быть записана в виде эквивалентного ей 


линейного операторного уравнения


\begin{equation}


Ax\equiv Ux+Bx = y \quad (x\in X,\ \ y\in Y),


\end{equation}


где $Ux=x^{(m)}(t)$, при краевых условиях (1.2). В силу условий теоремы 


операторы $A:X\to Y$ и $U:X\to Y$ имеют непрерывные обратные


$A^{-1}:Y\to X$ и $U^{-1}:Y\to X$.





Введем подпространства


\begin{gather*}


Y_n = \bigg\{\sum_{k=1}^n \beta_k t^{k-1}\bigg\} \equiv


\bH_{n-1},\quad \beta_k\in \R,\\


%


X_n = \{x_n\in \bH_{n+m-1}: R_l(x_n)=0,\ \ l=\ol{0,m-1}\}.


\end{gather*}


Ясно, что $X_n\subset X$, $Y_n\subset Y$ и $\dim X_n = \dim Y_n = n\in \N$.





Обозначим через $\Pi_n:Y\to Y_n$ оператор метода подобластей по любой из 


систем узлов (2.1) и (2.2). В силу следствий лемм 2 и 3 из [3], 


леммы 1.1 и ее следствия работы [2] для любой $f\in Y$ имеем


\begin{gather}


E_{n-1}(f)_Y\lea\|f-\Pi_n f\|_Y\lea\frac{\pi}{2}E_{n-1}(f)_Y,\quad n\in\N,\\


%


1\lea \|\Pi_n\|_{Y\to Y}\lea \frac{\pi}{2},\quad n\in\N.


\end{gather}





Запишем СЛАУ (1.4)--(1.5) в виде эквивалентного ей линейного операторного 


уравнения


\begin{equation}


A_n x_n\equiv Ux_n + \Pi_n B x_n = \Pi_n y \quad (x_n\in X_n, y_n\in Y_n)


\end{equation}


и покажем, что операторы $A_n:X_n\to Y_n$ при всех $n\gea n_0$ линейно 


обратимы, а обратные операторы ограничены по норме в совокупности:


\begin{equation}


\|A_n^{-1}\|_{Y_n\to X_n} = O(1),\quad n\to \infty.


\end{equation}


С этой целью к уравнениям (2.4) и (2.8) применим теорему 7 


из ([5], гл.\,1). Для любых $x_n\in X_n$, $x_n \ne 0$, из (2.4) и (2.8) находим


\begin{gather*}


\|Ax_n-A_nx_n\|_Y = \|B x_n - \Pi_n B x_n \|_Y =


\|x_n\|_X\|B z_n - \Pi_n B z_n\|_Y \lea \ve'_n\|x_n\|_X, \\


%


\ve'_n = \sup\begin{Sb} z_n\in X_n,\\ \|z_n\|_{X_n}=1\end{Sb}


\|B z_n - \Pi_n B z_n\|_Y\lea


\sup\begin{Sb} z\in X,\\ \|z\|_X\lea 1\end{Sb}


\|B z - \Pi_n B z\|_Y = \sup\limits_{f\in BS(0,1)}\|f-\Pi_n f\|_Y,


\end{gather*}


где $z_n = x_n/\|x_n\|$, $S(0,1)=\{x\in X:\|x\|_X\lea 1\}$.


Поскольку в силу условия~a) теоремы множество $BS(0,1)$ компактно 


в пространстве $Y$, то из (2.5) и теоремы Гельфанда (напр., [4],


с.\,274--276) следует


$\lim\limits_{n\to \infty} \ve'_n = 0$. Поэтому


\begin{equation}


\ve_n \equiv \|A-A_n\|_{X_n\to Y} \lea \ve'_n\to 0,\quad n\to \infty.


\end{equation}


Кроме того, в силу (2.5) для правых частей уравнений (2.4) и(2.7) имеем


\begin{equation}


\delta_n \equiv \|y-\Pi_n y\|_Y \lea


\frac{\pi}{2}E_{n-1}(y)_{L_2(\rho)}\to 0,\quad n\to\infty.


\end{equation}


Поэтому в силу теоремы 7 ([5], гл.\,1) для всех $n\in \N$ таких, что


\begin{equation}


q_n \equiv \ve_n \|A^{-1}\|_{Y\to X} <1 \quad (n>n_0),


\end{equation}


операторы $A_n:X_n\to Y_n$ непрерывно обратимы и для $A_n^{-1}:Y_n\to X_n$


справедливы оценки (2.8); кроме того, приближенные решения


$x_n^*=A_n^{-1}\Pi_n y$, определяемые по формуле (1.3$^*$), сходятся к 


точному решению $x^*=A^{-1}y$ в пространстве $X$ со скоростью


\begin{equation}


\|x^*-x^*_n\|_X = O(\ve_n + \delta_n),


\end{equation}


где $\ve_n$ и $\delta_n$ определены в (2.9) и (2.10).





Теперь с помощью теоремы 6 из ([5], гл.\,1) и соотношений (2.5)--(2.12) находим


\begin{gather}


\|x^*-x^*_n\|_X\lea \|E - A_n^{-1}\Pi_n A\|_{X\to X}


\|x^* - \widetilde{x}_n\|_X,\\


%


\|x^*-x^*_n\|_X \lea \|E- A_n^{-1}\Pi_n B\|_{X\to X} \|U^{-1}\|_{Y\to X}


\|Ux^*-\Pi_n U x^*\|_Y,


\end{gather}


где $\widetilde{x}_n$ --- произвольный элемент из $X_n$. Выбирая его 


соответствующим образом, из (2.13), (2.4), (2.6) получаем оценку


\begin{equation}


\|x^*-x^*_n\|_X = O\{E_{n+m-1}(x^*)_X\}.


\end{equation}


Аналогично, из (2.14), (2.4)--(2.6) и условия~б) теоремы находим


\begin{equation}


\|x^*-x^*_n\|_X = O\{E_{n-1}({x^*}^{(m)})_Y\}.


\end{equation}


Нетрудно показать, что


\begin{equation}


\|x^*-x^*_n\|_X \gea E_{n+m-1}(x^*)_X \gea E_{n-1}({x^*}^{(m)})_Y.


\end{equation}


Из неравенств (2.15)--(2.17) следуют соотношения (2.3).





\begin{thm}


В условиях теоремы $1$ приближенные решения $(1.3{}^*)$ сходятся в 


пространствах $C^k$ $(k=\ol{0,m-1}$, $m\gea 1)$ со скоростью


\begin{equation}


\|x^*-x_n^*\|_{C^k} = \suml_{i=0}^k


\|{x^*}^{(i)}(t) - {x_n^*}^{(i)}(t)\|_C =


O\{E_{n-1}({x^*}^{(m)})_{L_2(\rho)}\}.


\end{equation}


Если $y\in C$, а оператор $B:X\to C$ ограничен, то справедлива порядковая 


оценка


\begin{equation}


\|x^* - x^*_n\|_{C^m} = O\{E_{n-1}({x^*}^{(m)})_C \ln n\};


\end{equation}


если же $y\in C$ и


\begin{equation}


\|\Pi_n B\|_{X\to C} = O(1),\quad n\to \infty,


\end{equation}


то справедлива асимптотическая оценка


\begin{equation}


\|x^*-x^*_n\|_{C^m} \sim \|{x^*}^{(m)} - \Pi_n {x^*}^{(m)}\|_C,


\quad n\to \infty,


\end{equation}


где $E_{n-1}(\vp)_C$ --- наилучшее равномерное приближение функции


$\vp \in C$ алгебраическими многочленами степени не выше $n-1$ $(n\in \N)$,


а символ $\sim$ означает знак сильной эквивалентности.


\end{thm}





\begin{pf}


Оценка (2.18) является следствием теоремы 1 и того 


известного факта, что пространство $W_2^m(\rho)$ непрерывно вложено в 


любое из пространств $C^k$ при $k = 0,1,\dotsc, m-1$


($m\gea 1$, $C^0= C$). Переходя к доказательству (2.19), отметим, что


\begin{equation}


\|x^* - x^*_n\|_{C^m} = \|x^* - x^*_n\|_{C^{m-1}}+


\|{x^*}^{(m)}-{x_n^*}^{(m)}\|_C,


\end{equation}


где первое слагаемое оценено в (2.18), а для оценки второго слагаемого 


воспользуемся тождествами


\begin{equation}


{x^*}^{(m)} \equiv y(t) - B(x^*;t),\quad


{x^*_n}^{(m)} \equiv \Pi_n(y;t) - (\Pi_n B)(x^*_n;t),


\end{equation}


где $-1\lea t \lea 1$.





Из (2.23) следует тождество


\begin{equation}


{x^*}^{(m)}-{x^*_n}^{(m)} = ({x^*}^{(m)} -


\Pi_n {x^*}^{(m)}) - (\Pi_n B)(x^* - x^*_n).


\end{equation}


В силу теоремы 1 и ограниченности оператора $B:X\to C$ из (2.24) находим 


неравенства


\begin{multline}


\|{x^*}^{(m)}-{x_n^*}^{(m)}\|_C \lea \|{x^*}^{(m)}-\Pi_n{x^*}^{(m)}\|_C +


\|\Pi_n\|_{C\to C} \|B\|_{X\to C}\|x^* - x^*_n\|_X \lea \\


%


\lea 2 \|\Pi_n\|_{C\to C} E_{n-1}({x^*}^{(m)})_C + \|\Pi_n\|_{C\to C}


\|B\|_{X\to C}O\{E_{n-1}({x^*}^{(m)})_{L_2(\rho)}\} = \\


%


=\|\Pi_n\|_{C\to C}O\{E_{n-1}({x^*}^{(m)})_C +


E_{n-1}({x^*}^{(m)})_{L_2(\rho)}\} =


O\{\|\Pi\|_{C\to C} E_{n-1}({x^*}^{(m)})_C\}.


\end{multline}


Поскольку в силу (2.18)


\begin{equation}


\|x^* - x^*_n\|_{C^{m-1}} = O\{E_{n-1}({x^*}^{(m)})_{L_2(\rho)}\} =


O\{E_{n-1}({x^*}^{(m)})_C\},


\end{equation}


то с учетом (2.22), (2.25) и (2.26) имеем


\begin{equation}


\|x^*-x^*_n\|_{C^m} = O\{\|\Pi_n\|_{C\to C}E_{n-1}({x^*}^{(m)})_C\}.


\end{equation}


Известно [6], [7], [2], что для узлов (2.1) и (2.2) справедлива оценка


\begin{equation}


\|\Pi_n\|_{C\to C} = O(\ln n),\quad n\to \infty.


\end{equation}


Из (2.27) и (2.28) следует оценка (2.19).





Далее, из (2.24) с учетом (2.20) и теоремы 1 последовательно находим


\begin{gather*}


\|{x^*}^{(m)} - {x_n^*}^{(m)}\|_C \lea \|{x^*}^{(m)} -


\Pi_n {x^*}^{(m)}\|_C + \|\Pi_n B\|_{X\to C}


\|x^*-x^*_n\|_X \lea \\


%


\lea \|{x^*}^{(m)} - \Pi_n {x^*}^{(m)}\|_C +


O\{E_{n-1}({x^*}^{(m)})_{L_2(\rho)}\} \sim 


\|{x^*}^{(m)} - \Pi_n {x^*}^{(m)}\|_C,\quad n\to \infty; \\


%


\|{x^*}^{(m)} - {x^*_n}^{(m)}\|_C \gea \|{x^*}^{(m)} -


\Pi_n {x^*}^{(m)}\|_C - \|\Pi_n B\|_{X\to C}


\|x^* - x^*_n\|_X \sim \|{x^*}^{(m)}-\Pi_n {x^*}^{(m)}\|_C,\quad n\to \infty;


\end{gather*}


поэтому


\begin{equation}


\|{x^*}^{(m)} - {x^*_n}^{(m)}\|_C \sim


\|{x^*}^{(m)} - \Pi_n {x^*}^{(m)}\|_C,\quad n\to \infty.


\end{equation}


Из соотношений (2.22), (2.26), (2.28) и (2.29) следует оценка (2.21).


\end{pf}





Из теоремы 2 и прямых теорем теории равномерных приближений функций


(напр., [8]--[10]) получаем 





\begin{cornonum}


Пусть функция $y(t)$ и оператор $B$ таковы, что функция ${x^*}^{(m)}(t)$ 


удовлетворяет условию Дини--Липшица на сегменте $[-1,1]$. Тогда 


метод (1.1)--(1.5), (2.1)--(2.2) сходится в пространстве $C^m$ со 


скоростью (2.19); если же


\begin{equation}


{x^*}^{(m)}(t) \in W^r H^\alpha[-1,1]\quad (r+1\in \N,\ \ 0<\alpha\lea 1),


\end{equation}


то и со скоростью


\begin{equation}


\|x^* - x^*_n\|_{C^m} = O\bigg(\frac{\ln n}{n^{r+\alpha}}\bigg).


\end{equation}


\end{cornonum}








\section{Некоторые замечания и дополнения}





$1^\circ$. Отметим, что утверждения, аналогичные теореме 2 и ее следствию, 


могут быть получены также самостоятельно, но при более жестких, чем выше, 


условиях; напр., как и в ([5], гл.\,4), с использованием оценки (2.28),


доказывается





\begin{thm}


Пусть функции $y(t)$ и $B(x;t)$


равномерно\footnote{Это означает, что


$\sup\limits\{\omega(Bx;\delta):x\in C^m$,


$\|x\|_{C^m} = 1\} = o\big\{\frac{1}{|\ln \delta|}\big\}$,


$\delta \to +0$, где $\omega(\vp;\delta)$ --- модуль не\-пре\-рыв\-ности 


функции $\vp\in C[-1,1]$ с шагом $\delta\in (0,2]$.}


относительно $x\in C^m$ удовлетворяют на $[-1,1]$ условию Дини--Липшица.


Тогда СЛАУ $(1.4)$--$(1.5)$, $(2.1)$--$(2.2)$ однозначно разрешима при 


всех $n\gea n_1 \in \N$ и приближенные решения $(1.3{}^*)$ сходятся в 


пространстве $C^m$ со скоростью


$$


\|x^* - x^*_n\|_{C^m} \sim \|{x^*}^{(m)} - \Pi_n {x^*}^{(m)}\|_C =


O\{E_{n-1}({x^*}^{(m)})_C \ln n\},\quad n\to \infty;


$$


если же выполняется условие $(2.30)$, то и со скоростью $(2.31)$.


\end{thm}





$2^\circ$. С помощью результатов работы [11] и приведенных выше теорем 


доказывается





\begin{thm}


Метод осциллирующих функций $(1.1)$--$(1.5)$, $(2.1)$--$(2.2)$


обладает следующими оптимальными свойствами$:$


\begin{itemize}


\item[а)] в условиях любой из теорем $1$ и $2$ метод оптимален по порядку 


среди всевозможных полиномиальных методов решения задачи $(1.1)$--$(1.2);$


\item[б)] в условиях второй части теоремы $2$, а также теоремы $3$ метод 


является асимптотически оптимальным среди всевозможных методов 


осциллирующих функций, основанных на аппроксимации полиномами.


\end{itemize}


\end{thm}





$3^\circ$. Приведенные выше результаты остаются справедливыми и в том 


случае, когда в соболевском пространстве $X = W_2^m(\rho)$ вводится 


другая, но эквивалентная предыдущей, норма. Например, в силу условия~б)


теоремы 1 можно положить


\begin{equation}


\|x(t)\|_{W_2^m(\rho)} = \|x^{(m)}(t)\|_{L_2(\rho)} =


\|Ux\|_{L_2(\rho)},\quad x\in X;


\end{equation}


в некоторых случаях норма (3.1) удобнее, чем использованная выше. 


Это видно хотя бы из следующего простого результата.





\begin{thm}


Пусть уравнение $x^{(m)}(t)=0$ при краевых условиях $(1.2)$ имеет лишь 


тривиальное решение и выполнено неравенство


\begin{equation}


q = \frac{\pi}{2}\|B\|_{X\to Y} <1.


\end{equation}


Тогда как задача $(1.1)$--$(1.2)$, так и СЛАУ $(1.4)$--$(1.5)$,


$(2.1)$--$(2.2)$ для любых $n\in \N$ однозначно разрешимы при любых 


правых частях$;$ приближенные решения $(1.3{}^*)$ сходятся к точному


решению $x^*(t)$ в пространстве $X$ с нормой $(3.1)$,


при этом для любых $n\in\N$ справедлива оценка


\begin{equation}


\|x^* - x^*_n\|_{W^m_2(\rho)} \lea


\frac{\pi}{2-\pi\|B\|_{X\to Y}}E_{n-1}({x^*}^{(m)})_{L_2(\rho)}.


\end{equation}


\end{thm}





{\bf Доказательство. } Из (3.1) и первого условия теоремы следует, что 


оператор $U:X\to Y$ непрерывно обратим и


\begin{equation}


\|U\|_{X\to Y} = \|U^{-1}\|_{Y\to X} = 1.


\end{equation}


Поэтому уравнения (2.4) и (2.7) эквивалентны соответственно уравнениям


\begin{gather}


Kx \equiv x + U^{-1}Bx = U^{-1} y \quad (x,U^{-1}y\in X), \\


%


K_n x_n \equiv x_n + U^{-1}\Pi_n B x_n = U^{-1} \Pi_n y


\quad (x_n,U^{-1}\Pi_n y\in X_n);


\end{gather}


при этом операторы $A$ и $K$ (а также $A_n$ и $K_n$) обратимы


(или необратимы) одновременно и в случае их обратимости справедливы 


соотношения


\begin{equation}


\|A^{-1}\|_{Y\to X} = \|K^{-1}\|_{X\to X},


\quad \|A_n^{-1}\|_{Y_n\to X_n} = \|K_n^{-1}\|_{X_n\to X_n}.


\end{equation}


В силу (3.2), (3.4), (3.7) и полноты пространства $X$ с нормой (3.1) 


оператор $K:X\to X$ (следовательно, и $A:X\to Y$) непрерывно обратим и


\begin{equation}


\|A^{-1}\|_{Y\to X}\lea \frac{1}{1-\|U^{-1}B\|_{X\to X}} \lea \frac{\pi}{\pi-2}<\infty;


\end{equation}


аналогично, в силу (2.6), (3.2), (3.4), (3.7) и полноты пространства $X_n$ 


операторы $K_n:X_n\to X_n$ (следовательно, и $A_n:X_n\to Y_n$) непрерывно 


обратимы при любых $n\in \N$ и


\begin{equation}


\|A^{-1}_n\|_{Y_n\to X_n} \lea


\frac{1}{1 - \|U^{-1}\Pi_n B\|_{X\to X}}\lea \frac{1}{1-q}.


\end{equation}





В силу (3.5)--(3.9) первая часть теоремы доказана. Для доказательства 


второй части воспользуемся соотношениями (2.5), (2.6), (2.14), (3.1),


(3.2), (3.4), (3.8) и (3.9), из которых при любых $n\in\N$  находим


\begin{gather*}


\|x^* - x^*_n\|_X \lea (1+\|A^{-1}_n\|_{Y_n\to X_n}


\|\Pi_n\|_{Y\to Y_n} \|B\|_{X\to Y})\|U^{-1}\|_{Y\to X}


\|U x^* - \Pi_n Ux^*\|_Y \lea \\


%


\lea \bigg(1+\frac{1}{1-q}\,\frac{\pi}{2}\|B\|_{X\to Y}\bigg)


\frac{\pi}{2}E_{n-1}(Ux^*)_Y \lea


\frac{1}{1-q}\frac{\pi}{2}E_{n-1}(Ux^*)_Y \lea


\frac{\pi}{2-\pi\|B\|_{X\to Y}}E_{n-1}({x^*}^{(m)})_{Y}.\quad\square


\end{gather*}








\begin{thebibliography}{99}





\bibitem{1}


Ермолаева Л.Б. {\em Решение линейных уравнений методом осциллирующих


функций} // Тр. Матем. центра им.~Н.И.\,Лобачевского. Актуальные проблемы 


математики и механики. -- Казань: Унипресс, 2000. -- Т.\,5. -- С.\,83--84.


\bibitem{2}


Ермолаева Л.Б. {\em Решение интегральных уравнений методом подобластей}


// Изв. вузов. Математика. -- 2002. -- \No\,9. -- С.\,37--49.


\bibitem{3}


Габдулхаев Б.Г., Ермолаева Л.Б. {\em Интерполяционные полиномы Лагранжа 


в пространствах Соболева} // Изв. вузов. Математика. -- 1997.


-- \No\,5. -- С.\,7--19.


\bibitem{4}


Канторович Л.В., Акилов Г.П. {\em Функциональный анализ в нормированных


пространствах}. -- М.: Физматгиз, 1959. -- 684~с.


\bibitem{5}


Габдулхаев Б.Г. {\em Оптимальные аппроксимации решений линейных задач}.


-- Казань: Изд-во Казанск. ун-та, 1980. -- 232~с.


\bibitem{6}


Нагих В.В. {\em Оценка нормы некоторого полиномиального оператора в


пространстве непрерывных функций} // Методы вычислений. Вып.\,10.


-- Л.: Изд-во ЛГУ, 1976. -- С.\,99--102.


\bibitem{7}


Керге Р.М. {\em О сходимости и устойчивости метода подобластей}:


Автореф. дис.~$\dotsc$ канд. физ.-матем. наук. -- Тарту, 1979. -- 10~с.


\bibitem{8}


Тиман А.Ф. {\em Теория приближения функций действительного переменного}.


-- М.: Физматгиз, 1960. -- 624~с.


\bibitem{9}


Даугавет И.К. {\em Введение в теорию приближения функций}. -- Л.: 


Изд-во ЛГУ, 1976. -- 184~с.


\bibitem{10}


Дзядык В.К. {\em Введение в теорию равномерного приближения функций 


полиномами}. -- М.: Наука, 1977. -- 490~с.


\bibitem{11}


Габдулхаев Б.Г. {\em Оптимизация прямых и проекционных методов решения 


операторных уравнений} // Изв. вузов. Математика.


-- 1999. -- \No\,12. -- С.\,3--18.





\end{thebibliography}





\vskip 2cm





\post{\it Казанская государственная}{\it Поступила}


\post{\it архитектурно-строительная академия}{26.06.2002}





\label{end}


\end{document}


